\section{Introduction}
\label{sec:introduction}

Scientific conclusions depend on more than carrying out a numerical operation correctly. Consider an adsorbate that moves across a catalyst surface without appreciably changing its internal geometry. Its displacement relative to the surface and its change in internal structure can both be reported in angstroms, but they answer different questions. A language model can select a genuine measurement, perform the requested arithmetic correctly, and still make an unsupported claim because the measurement describes the wrong physical quantity. The difficulty is not simply numerical accuracy: it is whether the \emph{definition} of the chosen evidence licenses the inference.

Tool-using language models make this distinction important. ChemCrow links language models to chemistry tools, and program-aided language models execute model-generated calculations rather than requiring the language model to perform all arithmetic itself.\cite{bran2024augmenting,gao2023pal} These approaches demonstrate the usefulness of external tools, but correct execution does not by itself establish that the model selected a physically compatible input. Converting units, for example, cannot turn a measure of molecular translation into a measure of internal deformation. Here we ask whether checking the numerical calculation and the source comparison is sufficient, or whether a checker must also validate the measurement's physical meaning.

Scientific evidence evaluation is not a new problem. SciBench evaluates scientific problem solving across mathematics, chemistry, and physics; SciCUEval studies information selection, missing-information detection, and context-aware inference in scientific materials.\cite{wang2024scibench,yu2026scicueval} ALCE separately assesses generated answers and the support provided by citations, while GSM-Symbolic uses controlled variants to probe sensitivity to problem presentation.\cite{gao2023enabling,mirzadeh2025gsm} GeoScope does not claim to introduce evidence selection or robustness testing. It isolates a narrower, physically specified question: whether a selected atomistic measurement and an allowed mathematical relation really prove a claim under declared coordinate conventions.

Atomistic calculations offer a suitable test because different measurements can be constructed from the same structure pair. The Open Catalyst 2020 dataset and OC20-Dense supply adsorbate--surface structures, and UMA provides a pretrained potential for producing relaxed coordinate endpoints.\cite{chanussot2021open,lan2023adsorbml,wood2025uma} We distinguish mobile-atom root-mean-square deviation (RMSD), adsorbate RMSD in the stored coordinate frame, and internal adsorbate RMSD after periodic unwrapping and rigid alignment. Each uses a different combination of atom selection and coordinate handling. When definitions agree, an adsorbate subset can support a valid bound on mobile-atom RMSD; when definitions do not agree, two plausible-looking distances need not support the same comparison. The definitions and admissible operations are supplied to the models rather than left for them to discover.

We define a \emph{licensed proof} as an evidence record and operation whose physical definitions permit the claimed conclusion. A correct verdict with an incompatible record is not a valid proof. Conversely, a missing exact measurement need not imply insufficient evidence when an allowed one-sided bound settles the claim. These distinctions let us compare progressively stricter checking of the \emph{same} raw model selection: arithmetic-only checking, comparison-ID checking, and full-definition checking. A rejection by the checker is not the same as a scientifically justified conclusion of insufficient evidence. Nor does rejection find replacement evidence.

The original GeoScope study contains 1,760 claim variants from 55 qualified source groups, represented as structured records and short reports. Three compact open-weight models answer directly or select an evidence record and operation. A deterministic oracle tests whether the supplied records and declared rules can resolve each case. We subsequently evaluate four models on 48 independently sourced OC20 IS2RE/IS2RS trajectories and measure decision-token confidence on the frozen external benchmark. This later study tests recurrence of the observed failure mechanisms beyond the parent ALCHEMI/OC20-Dense structures, but does not establish unseen model-pretraining data or independent physical ground truth.

Our contributions are:
\begin{itemize}
  \item A reproducible atomistic evidence benchmark that makes atom sets, coordinate frames, periodic handling, and permitted proof operations explicit, including resolvable subset bounds and genuinely insufficient evidence.
  \item A matched-selection comparison showing which invalid numerical acceptances survive arithmetic and comparison-ID checks but are excluded by full physical-definition validation.
  \item A failure-inclusive evaluation reporting accepted-proof correctness and \emph{valid numerical coverage} separately, together with evidence-presentation and operation-specific diagnostics and a deterministic oracle baseline.
  \item A four-model external-source evaluation and matched confidence replay that distinguish strict output-format compliance, physical-proof correctness, confidence ranking, and scientifically justified abstention.
\end{itemize}

Our primary matched-selection result is that full-definition checking removes 81 additional invalid acceptances for Qwen while preserving all 312 licensed numerical answers. Those answers still cover only 28.36\% of the original 1,100 answerable cases. The external and confidence analyses ask whether this interpretation generalizes across source collections and whether model-internal token scores can identify valid proofs. Throughout, rejecting unsupported evidence, recovering usable evidence, and knowing when evidence is insufficient remain separate tasks.
